\documentclass[11pt,a4paper]{article}
\usepackage[utf8]{inputenc}
\usepackage[T1]{fontenc}
% Scommentare in locale se la texlive li include:
%\usepackage[italian]{babel}
%\usepackage{lmodern}
%\usepackage{microtype}
\usepackage[margin=2.6cm]{geometry}
\usepackage{booktabs}
\usepackage{amsmath}
\usepackage{amssymb}
\usepackage{array}
\usepackage{xcolor}
\usepackage{listings}
\usepackage{hyperref}
\usepackage{enumitem}
\hypersetup{colorlinks=true, linkcolor=blue!50!black, urlcolor=blue!50!black}

\lstset{
  basicstyle=\ttfamily\small, breaklines=true, frame=single, framerule=0.2pt,
  backgroundcolor=\color{gray!6}, showstringspaces=false, language=Python
}

\newcommand{\code}[1]{\texttt{#1}}
\newcommand{\flow}[1]{\texttt{\small #1}}

\title{\textbf{Pipeline SDMX--ISTAT per popolazioni sintetiche comunali e sub-comunali}\\[2mm]
\large Note tecniche --- v0.3}
\author{M.~Degli Esposti}
\date{19 luglio 2026 \quad (v0.1--v0.2: 18--19 luglio 2026)}

\begin{document}
\maketitle

\begin{abstract}
\noindent
Pipeline riproducibile dai dati pubblici ISTAT alla popolazione sintetica
MaxEnt, ora estesa al livello \emph{sub-comunale}: Brescia K7C con variabile
zona a 33 quartieri ($|X|=177\,408$, $m\simeq2\,500$), da sezioni di
censimento 2023. Risultati principali aggiunti in v0.3:
(vi) le sezioni di censimento portano i codici delle aree sub-comunali
(\code{COM\_ASC1}) per tutti i comuni partizionati --- niente spatial join;
(vii) quinta identit\`a esatta della base riconciliata: sezioni $\equiv$
censimento comunale 2023 $\equiv$ anagrafe 1/1/2024 $=198\,259$ alla persona;
(viii) \textbf{lezione di consistenza generalizzata}: con vincoli di
uguaglianza devono coincidere \emph{tutti} i margini impliciti sovrapposti tra
blocchi, non solo quelli verso la spine --- violarla produce collasso del
duale; l'IPF entra come \emph{armonizzatore di tabelle} a monte del MaxEnt
(non come solver); (ix) terza verifica della formula del pavimento stocastico
PCD (predetto $0.0650$, osservato $0.0634$, $N=2\cdot10^5$): da osservazione a
risultato replicato; (x) il pavimento limita lo \emph{stimatore-pool}, non la
distribuzione appresa: la MRE dei $\lambda$ mediati da Adam batte il pavimento
di $\sim2\times$.
\end{abstract}

\tableofcontents

%=======================================================================
\section{Scopo e contesto}
Obiettivo: dato un codice comune ISTAT, produrre il miglior constraint set dai
dati pubblici e la popolazione MaxEnt corrispondente; caso guida Brescia (SIN
Caffaro), ora con attributo di zona a 33 quartieri. Esecuzione locale (WSL2,
env \code{ml}); solver dal repo \code{maxent-popsynth-pcd} (esatto L-BFGS,
GibbsPCD con kernel Numba CSR).

%=======================================================================
\section{Endpoint SDMX \code{esploradati}: comportamenti non documentati}
\label{sec:quirks}
(v0.1, invariata) \begin{enumerate}[leftmargin=*,itemsep=1pt]
  \item limite $\sim$260 char/segmento URL (http.sys) $\Rightarrow$ wildcard;
  \item \code{startPeriod}/\code{endPeriod} ignorati $\Rightarrow$ filtro
  client-side; serie storiche complete gratis;
  \item SDMX-CSV via \code{Accept: application/vnd.sdmx.data+csv;version=1.0.0};
  \item rate limit $\sim$5/min, throttling 13\,s, retry;
  \item stime censuarie comunali \emph{non arrotondate} (decimali, additive).
\end{enumerate}

%=======================================================================
\section{Le fonti: tre laghi pi\`u uno}
\label{sec:laghi}
Ai tre laghi di v0.1 (anagrafe \flow{DCIS\_*}; censimento permanente
\flow{DF\_DCSS\_*}; indagini campionarie, solo regionali) si aggiunge il
canale \textbf{sub-comunale}:
\begin{itemize}[leftmargin=*,itemsep=1pt]
  \item le \emph{aree sub-comunali} (ASC) aggregate via SDMX
  (\flow{DF\_BULK\_SUBCOM2021/2023}) coprono \textbf{solo i capoluoghi delle
  Citt\`a Metropolitane}: Brescia esclusa;
  \item le \emph{sezioni di censimento} coprono tutti i comuni: bulk download
  (niente SDMX, niente rate limit)
  \code{Dati\_regionali\_2023.zip} su
  \code{esploradati.istat.it/databrowser/DWL/PERMPOP/SUBCOM/} --- un xlsx per
  regione + \code{TRACCIATO FILE REGIONALI.xlsx}; edizione censimento
  31-12-2023 (agg.\ 5/3/2026: pi\`u variabili, +59 comuni); disponibile anche
  il 2021;
  \item \textbf{scoperta chiave}: il file sezioni contiene le colonne
  \code{COM\_ASC1/2/3} --- i codici delle aree sub-comunali sono assegnati
  sezione per sezione \emph{per tutti i comuni partizionati}. Per Brescia
  \code{COM\_ASC1} ha 33 valori (\code{17029001}--\code{17029033}) $=$ i 33
  quartieri: lo spatial join sezioni$\to$quartieri \`e superfluo, basta un
  groupby.
\end{itemize}
Variabili del tracciato utili: sesso$\times$et\`a quinquennale (P30--45,
P67--82), istruzione a 5 livelli per sesso (P86--100, universo 9+), occupati
15--64 per sesso (P101--103), italiani/stranieri per sesso$\times$macro-et\`a
(IT1--12, ST22--33), UE/extra-UE, background migratorio (EM1--6), famiglie
(PF), abitazioni.

%=======================================================================
\section{Brescia: sezioni, quartieri, audit}
\label{sec:submun}
1\,822 sezioni, 33 quartieri; 1 sezione fittizia (888888x, senza fissa
dimora): 549 persone, assegnate da ISTAT al quartiere \code{17029004} ---
\textbf{tenute} (opzione a: coerenza contabile con i totali ufficiali;
esclusione con riscalatura come variante documentata per l'analisi di
esposizione Caffaro).

\medskip\noindent\textbf{Quinta identit\`a esatta della base riconciliata:}
\[
\text{sezioni 2023} \;\equiv\; \text{censimento comunale 2023} \;\equiv\;
\text{anagrafe 1/1/2024} \;=\; 198\,259 \ \text{(alla persona; stranieri
} 37\,478\text{)}.
\]
Additivit\`a interna piena (classi et\`a $\to$ P1; istruzione $\to$ P83).
Differenza di convenzione: le sezioni sono in \emph{interi} (allocazione
RBI$\leftrightarrow$RSBL di individui interi), il canale SDMX comunale in
stime decimali --- scarti di unit\`a possibili sui sotto-totali, assorbiti dal
disegno condizionale (\S\ref{sec:k7c}).

\medskip\noindent Ancoraggio temporale K7C: \code{--anno 2024}
(anagrafe 1/1/2024 $\equiv$ censimento 31-12-2023 $\equiv$ sezioni), tutto
sulla stessa data. Tabelle di quartiere (\code{build\_zona\_tables.py}):
Z1 zona$\times$sesso$\times$et\`a5 (1\,056 righe), Z2
zona$\times$sesso$\times$macro-et\`a$\times$ITL/FRG (396), Z3
zona$\times$sesso$\times$istruzione5 (330), Z4
zona$\times$sesso$\times$occupato/non (132, 15--64). Sanity qualitativa:
top-5 quota stranieri = Fiumicello (0.371), Centro Storico Nord (0.303),
Primo Maggio (0.285), Don Bosco (0.276), Chiesanuova (0.259) --- la geografia
sociale dell'orbita SIN Caffaro.

%=======================================================================
\section{K7C: disegno dei blocchi di zona}
\label{sec:k7c}
Livello \textbf{K7C}: \code{zona} in testa a \code{VAR\_ORDER};
$|X| = 33\times5\,376 = 177\,408$. Scelta di ariet\`a: \emph{piena} dove il
dato la sostiene (Z2 a ariet\`a 4: zona$\times$sesso$\times$et\`a$\times$
cittadinanza cattura direttamente ``dentro Fiumicello gli stranieri sono
giovani'', non solo i due fatti separati).

\medskip\noindent\textbf{Principio}: ogni tabella di zona entra come
$P(\text{zona}\mid\text{gruppo})\times$ conteggi comunali del gruppo, mai come
conteggi grezzi: i margini sommati sulle zone coincidono coi blocchi comunali
per costruzione. Assunzioni aggiuntive dichiarate: (2) quota di zona costante
entro la classe quinquennale (Z1) e la macro-classe (Z2); (3)
$P(\text{zona}\mid\text{sesso}, \text{occupato})$ costante sui bin 15--64
(Z4, che vincola il solo lato occupato).

\subsection{Lezione di consistenza generalizzata (e collasso istruttivo)}
\label{sec:lesson}
Il primo fit K7C \`e collassato con firma nuova: discesa (MRE
$0.26\to0.06$ in 40 iter), poi \emph{risalita} fino a plateau ($3.4$).
Causa: i blocchi Z definivano gli stessi pattern impliciti attraverso
decomposizioni condizionali diverse --- es.\ sul bin 15--24 il totale
(zona, sesso, et\`a) secondo Z1 usa le quote quinquennali 15--19/20--24,
secondo Z2 (sommato sulle cittadinanze) la quota piatta della macro-classe
15--64: due valori diversi per lo stesso pattern $\Rightarrow$ vincoli di
uguaglianza contraddittori $\Rightarrow$ duale illimitato.

\medskip\noindent\textbf{Regola generalizzata} (estende la lezione E/F di
v0.2): \emph{con vincoli di uguaglianza devono coincidere numericamente
TUTTI i margini impliciti sovrapposti tra blocchi --- non solo quelli verso
la spine; l'audit deve verificarli tutti prima del fit.}

\medskip\noindent\textbf{Fix: IPF come armonizzatore di tabelle} (non come
solver). Per ogni (sesso, bin): IPF sulla matrice 33$\times$2
zona$\times$cittadinanza con margini-riga $=$ Z1 e margini-colonna $=$ blocco
comunale B (Z2); per sesso: zona$\times$istruzione con righe Z1 e colonne
totali comunali (Z3); Z4 con cap cellwise su Z1 e ridistribuzione (34 celle
cappate su Brescia). Audit post-armonizzazione: Z2-vs-Z1 e Z3-vs-Z1
$\max|\Delta| \sim 10^{-11}$; Z-vs-comunale $=0$. Ruolo concettuale: il
raking, perdente come solver (AISTAT), \`e lo strumento giusto per rendere
mutuamente consistenti tabelle sovrapposte \emph{a monte} del MaxEnt.

%=======================================================================
\section{Fit K7C}
\label{sec:fitk7c}
Setup: \code{fit\_cs.py} (generalizza \code{fit\_k6c.py}; firme dei blocchi
per nome variabile, robuste allo shift di indici della zona), pruning
\code{--min-alpha 2e-4}, $m=2\,216$ effettivi su $2\,513$; esclusioni
$\alpha=0$ fuori dai solver, post-hoc sul supporto (5\,544 celle:
$2\times33\times2\times6\times7$). Matrice indicatrice $F$: 3.15\,GB,
costruita in 4\,s.

\begin{center}
\small
\begin{tabular}{@{}lrrrrr@{}}
\toprule
 & \textbf{MRE}$_{\alpha>0}$ & \textbf{H (nat)} & \textbf{supporto} & \textbf{KL vs exact} & \textbf{tempo}\\
\midrule
Esatto (L-BFGS) & $3.6\cdot10^{-4}$ & 9.157 & 130\,161/177\,408 & --- & 351\,s\\
GibbsPCD ($N{=}2\cdot10^5$, numba) & $3.1\cdot10^{-2}$ & 9.592 & 171\,864 & $1.2\cdot10^{-1}$ & 1\,070\,s\\
\bottomrule
\end{tabular}
\end{center}
\noindent Nota su H: se la zona fosse indipendente dagli attributi,
$H \approx H_{K6C} + H(\text{zona}) \approx 5.84 + 3.4$; il deficit misura
l'informazione mutua zona$\leftrightarrow$attributi catturata dal CS.
Campione $n=198\,259$ dal modello esatto (\code{popolazione\_K7C.csv}, con
colonna quartiere); marginali di zona e quote stranieri per quartiere
riprodotti (check sostanziale superato).

\subsection{Pavimento stocastico: terza verifica (risultato replicato)}
\[
\text{MRE}_{\min} \approx \frac{1}{m}\sum_j
\sqrt{\frac{1-\alpha_j}{\alpha_j N}}:
\qquad
\begin{array}{lll}
\text{K6C},\ N{=}5\cdot10^4: & \text{pred.\ }0.0597, & \text{oss.\ }0.051/0.074;\\
\text{K7C},\ N{=}2\cdot10^5: & \text{pred.\ }0.0650, & \text{oss.\ }0.0634.
\end{array}
\]
Su un CS $9\times$ pi\`u grande e con spettro di $\alpha$ diverso, la
predizione a priori sbaglia del 2\%: il pavimento \`e calcolabile dal solo
constraint set, prima di ogni run. Regole di dimensionamento:
$N \gtrsim 100/\alpha_{\min}$; per MRE target $\tau$:
$N \approx \overline{1/\alpha}/\tau^2$.

\subsection{Stimatore-pool vs distribuzione appresa (finding)}
\label{sec:estimator}
A convergenza, MRE calcolata sulle \emph{frequenze empiriche del pool}
$=0.063$ (al pavimento); MRE calcolata sulla \emph{distribuzione implicata
dai $\lambda$} (via matrice indicatrice) $=0.031$: fattore $\sim2$. I
$\lambda$ sono mediati nel tempo da Adam lungo centinaia di outer e integrano
l'informazione di molti pool successivi: \textbf{il pavimento stocastico
limita lo stimatore-pool (singolo snapshot), non la distribuzione appresa}.
Corollario: se il deliverable \`e $p_\lambda$ (o un campione da essa) anzich\'e
il pool finale, la precisione effettiva \`e migliore del pavimento nominale.

\subsection{Decisione di produzione (ribadita)}
Anche a $|X|=1.8\cdot10^5$ il solver esatto domina: $100\times$ pi\`u preciso
a un terzo del tempo. GibbsPCD: strumento per domini non enumerabili,
con $N$ dimensionato dalla formula; ruolo attuale, caratterizzazione dello
strumento (``controllability first'').

%=======================================================================
\section{Architettura aggiornata della pipeline}
\begin{center}
\begin{tabular}{@{}lp{9.4cm}@{}}
\toprule
\textbf{Script} & \textbf{Funzione}\\
\midrule
\code{istat\_catalog.py} & catalogo 4\,884 dataflow + grep\\
\code{istat\_sdmx.py} & modulo SDMX generico (KEY dal DSD, CSV, cache, retry)\\
\code{fetch\_comune.py} & rosa dei fondamentali comunali per codice comune\\
\code{build\_constraints.py} & tavole comunali $\to$ vincoli C1--C6 + nazionalit\`a\\
\code{build\_zona\_tables.py} & sezioni $\to$ tabelle di quartiere Z1--Z4 + audit\\
\code{cs\_build.py} (v2) & vincoli $\to$ ConstraintSet; livelli K6C/K7C; IPF di armonizzazione\\
\code{fit\_cs.py} & fit esatto + GibbsPCD (K6C/K7C), diagnostiche, popolazione\\
\bottomrule
\end{tabular}
\end{center}

%=======================================================================
\section{Sintesi dei livelli}
\begin{center}
\small
\begin{tabular}{@{}lrrrrr@{}}
\toprule
 & \textbf{anno} & $|X|$ & $m$ & \textbf{MRE esatto} & \textbf{H (nat)}\\
\midrule
K6C & 1/1/2025 & 5\,376 & 269 & $5.0\cdot10^{-4}$ & 5.842\\
K7C & 1/1/2024 & 177\,408 & 2\,513 & $3.6\cdot10^{-4}$ & 9.157\\
\bottomrule
\end{tabular}
\end{center}

%=======================================================================
\section{Da fare / prossimi passi}
\begin{itemize}[leftmargin=*,itemsep=1pt]
  \item[$\boxtimes$] Livello sub-comunale K7C (sezioni $\to$ 33 quartieri,
  IPF, fit) --- v0.3.
  \item[$\square$] \textbf{Run B} (pool $=$ popolazione $+$ $\lambda$ caldi
  dall'esatto): misura diretta della fluttuazione stazionaria $=$ pavimento
  isolato dal transiente; richiede warm start nel GibbsPCDSolver (subclass).
  \item[$\square$] Nazionalit\`a two-stage sul campione K7C
  (\code{nationality\_conditional.csv}); poi via/civico entro quartiere
  (macchinario v3).
  \item[$\square$] Seed esplicito nel kernel Numba (riproducibilit\`a tra
  backend e tra run).
  \item[$\square$] Variante Caffaro: esclusione sezione fittizia (549) con
  riscalatura, per l'analisi di esposizione.
  \item[$\square$] $\sigma$ campionarie del censimento permanente; secondo
  comune (Torino 001272); tavole aggiuntive (famiglie, background
  migratorio, UE/extra-UE come raffinamento cittadinanza).
  \item[$\square$] Aggiornamento periodico: sub-comunale 2024 quando
  rilasciato (piano ISTAT: annuale) $\to$ riallineamento in avanti.
\end{itemize}

\end{document}
